% solution source: 2011_NIGHT_sol.pdf page 1 (Night Observation Q1 -- Little
% Dolphin, marking scheme). The official solution is the marking scheme; no
% model drawings are given.
\textbf{On Drawing 1:}

Draw the view of the constellation Delfinus (Del) through the finder
scope: % sic -- "Delfinus" for Delphinus in the source
\begin{itemize}
\item all 4 stars of the parallelogram \hfill(1 point)
\item the `tail' \hfill(1 point)
\item correct scale \hfill(1 point)
\item correct shape \hfill(1 point)
\item additional field stars \hfill(1 point)
\end{itemize}

With an arrow, mark the apparent direction of motion of the stars across the
field of view of the finder scope caused by the rotation of the Earth.
\hfill(2 points)

Label the stars with the Bayer designations given on the map ($\alpha$,
$\beta$, $\gamma$, $\delta$ and $\epsilon$). \hfill(1 point)

Also label the brightest of these 5 stars ``$m_{\max}$''. \hfill(1 point)

Also label the faintest of these 5 stars ``$m_{\min}$''. \hfill(1 point)

\textbf{On Drawing 2:}

Draw the view of the Little Dolphin through the larger telescope:
\begin{itemize}
\item all 4 stars of the parallelogram \hfill(4 points)
\item the `tail' \hfill(1 point)
\item correct scale \hfill(1 point)
\item correct shape \hfill(1 point)
\item additional field stars \hfill(1 point)
\end{itemize}

With an arrow, mark the apparent direction of motion of the stars across the
field of view of the finder scope caused by the rotation of the Earth.
% sic -- the source says "finder scope" here although Drawing 2 is the view
% through the larger telescope
\hfill(2 points)

Label the stars of the Little Dolphin $\alpha'$, $\beta'$, $\gamma'$,
$\delta'$ and $\epsilon'$ such that they match the labels of the stars in
the constellation Delphinus as given on the map. \hfill(3 points)

Label the brightest of these 5 stars ``$m_{\max}$''. \hfill(2 points)

Total: 25 points.
